% GENERATED FILE. Edit canonical lessons, then run npm run generate:books.
% canonical-source-hash: fnv1a64:d72f881627b33fcf
% canonical-lessons: UR-C134-gila, UR-C134-sukha, UR-C134-chaura, UR-C134-tang, UR-C134-gahra

\chapter{Wet, Dry, Wide, Narrow, Deep}
\label{ch:ur-wet-dry-wide-narrow-deep}
\hlchaptermodality{134}

\begin{chapteropening}
\textbf{By the end of this chapter:} \emph{I can say wet, dry, wide, narrow and deep.}

Complete the last lesson of chapter 134: I can say wet, dry, wide, narrow and deep.
\end{chapteropening}

\section[\texorpdfstring{\ur{گیلا} (gīlā)}{گیلا (gīlā)}]{\ur{گیلا} (gīlā) — wet}

\label{lesson:UR-C134-gila}



\begin{quote}
Before the new one: say the Urdu for soft, then the Urdu for hard.
\end{quote}

\subsection*{You'll want to know: \ur{گیلا}}
\textbf{\ur{گیلا}} — \emph{gīlā} — \textquotedblleft{}wet\textquotedblright{}.

\begin{cousinweb}
\textbf{\ur{گیلا}} (\emph{gīlā}) is a tadbhava, a word inherited and reshaped by sound change. It comes from \emph{*grilla}, \textquotedblleft{}wet, damp\textquotedblright{}, a form CDIAL reconstructs from its descendants (CDIAL 4386). CDIAL calls a further link with \emph{*gṛdla} very doubtful.
\end{cousinweb}

\begin{grammarlens}[title={What you've built}]
One more word for this chapter.
\end{grammarlens}

\subsection*{Your turn}
\begin{itemize}
\raggedright
  \item \emph{Say it:} \emph{gīlā}
  \item \emph{Say it:} \emph{gīlā}, once more
  \item \emph{Say it:} say \emph{gīlā}
  \item \emph{Recall:} say the Urdu for soft, then the Urdu for hard, then say \emph{gīlā} again
\end{itemize}

\begin{tcolorbox}[breakable,colback=teal!4,colframe=teal!35!black,title={Before you move on}]
What does \ur{گیلا} mean? (\textquotedblleft{}Wet\textquotedblright{}.) Say it once more.
\end{tcolorbox}
\section[\texorpdfstring{\ur{سوکھا} (sūkhā)}{سوکھا (sūkhā)}]{\ur{سوکھا} (sūkhā) — dry}

\label{lesson:UR-C134-sukha}



\begin{quote}
Before the new one: say the Urdu for hard, then the Urdu for wet.
\end{quote}

\subsection*{You'll want to know: \ur{سوکھا}}
\textbf{\ur{سوکھا}} — \emph{sūkhā} — \textquotedblleft{}dry\textquotedblright{}.

\begin{cousinweb}
\textbf{\ur{سوکھا}} (\emph{sūkhā}) is a tadbhava, a word inherited from Sanskrit and reshaped by sound change. It comes from \emph{śuṣka}, \textquotedblleft{}dried\textquotedblright{} (CDIAL 12548).
\end{cousinweb}

\begin{grammarlens}[title={What you've built}]
One more word for this chapter.
\end{grammarlens}

\subsection*{Your turn}
\begin{itemize}
\raggedright
  \item \emph{Say it:} \emph{sūkhā}
  \item \emph{Say it:} \emph{sūkhā}, once more
  \item \emph{Say it:} say \emph{sūkhā}
  \item \emph{Recall:} say the Urdu for hard, then the Urdu for wet, then say \emph{sūkhā} again
\end{itemize}

\begin{tcolorbox}[breakable,colback=teal!4,colframe=teal!35!black,title={Before you move on}]
What does \ur{سوکھا} mean? (\textquotedblleft{}Dry\textquotedblright{}.) Say it once more.
\end{tcolorbox}
\section[\texorpdfstring{\ur{چوڑا} (chauṛā)}{چوڑا (chauṛā)}]{\ur{چوڑا} (chauṛā) — wide}

\label{lesson:UR-C134-chaura}



\begin{quote}
Before the new one: say the Urdu for wet, then the Urdu for dry.
\end{quote}

\subsection*{You'll want to know: \ur{چوڑا}}
\textbf{\ur{چوڑا}} — \emph{chauṛā} — \textquotedblleft{}wide\textquotedblright{}.

\begin{cousinweb}
\textbf{\ur{چوڑا}} (\emph{chauṛā}) is a tadbhava, a word inherited and reshaped by sound change. It comes from \emph{*ca-uḍa}, \textquotedblleft{}wide, open\textquotedblright{}, a form CDIAL reconstructs from its descendants (CDIAL 4534). CDIAL suggests, with a \textquotedblleft{}possibly\textquotedblright{}, that it was built on \emph{cau-}, \textquotedblleft{}four\textquotedblright{}.
\end{cousinweb}

\begin{grammarlens}[title={What you've built}]
One more word for this chapter.
\end{grammarlens}

\subsection*{Your turn}
\begin{itemize}
\raggedright
  \item \emph{Say it:} \emph{chauṛā}
  \item \emph{Say it:} \emph{chauṛā}, once more
  \item \emph{Say it:} say \emph{chauṛā}
  \item \emph{Recall:} say the Urdu for wet, then the Urdu for dry, then say \emph{chauṛā} again
\end{itemize}

\begin{tcolorbox}[breakable,colback=teal!4,colframe=teal!35!black,title={Before you move on}]
What does \ur{چوڑا} mean? (\textquotedblleft{}Wide\textquotedblright{}.) Say it once more.
\end{tcolorbox}
\section[\texorpdfstring{\ur{تنگ} (tang)}{تنگ (tang)}]{\ur{تنگ} (tang) — narrow, tight}

\label{lesson:UR-C134-tang}



\begin{quote}
Before the new one: say the Urdu for dry, then the Urdu for wide.
\end{quote}

\subsection*{You'll want to know: \ur{تنگ}}
\textbf{\ur{تنگ}} — \emph{tang} — \textquotedblleft{}narrow, tight\textquotedblright{}.

\begin{cousinweb}
\textbf{\ur{تنگ}} (\emph{tang}) is a loan, usually traced to Persian \emph{tang}, \textquotedblleft{}narrow, tight\textquotedblright{}.
\end{cousinweb}

\begin{grammarlens}[title={What you've built}]
One more word for this chapter.
\end{grammarlens}

\subsection*{Your turn}
\begin{itemize}
\raggedright
  \item \emph{Say it:} \emph{tang}
  \item \emph{Say it:} \emph{tang}, once more
  \item \emph{Say it:} say \emph{tang}
  \item \emph{Recall:} say the Urdu for dry, then the Urdu for wide, then say \emph{tang} again
\end{itemize}

\begin{tcolorbox}[breakable,colback=teal!4,colframe=teal!35!black,title={Before you move on}]
What does \ur{تنگ} mean? (\textquotedblleft{}Narrow, tight\textquotedblright{}.) Say it once more.
\end{tcolorbox}
\section[\texorpdfstring{\ur{گہرا} (gahrā)}{گہرا (gahrā)}]{\ur{گہرا} (gahrā) — deep, dark (of colour)}

\label{lesson:UR-C134-gahra}



\begin{quote}
Before the new one: say the Urdu for wide, then the Urdu for narrow.
\end{quote}

\subsection*{You'll want to know: \ur{گہرا}}
\textbf{\ur{گہرا}} — \emph{gahrā} — \textquotedblleft{}deep, dark (of colour)\textquotedblright{}.

\begin{cousinweb}
\textbf{\ur{گہرا}} (\emph{gahrā}) is a tadbhava, a word inherited from Sanskrit and reshaped by sound change. It comes from \emph{gabhīra}, \textquotedblleft{}deep, solemn\textquotedblright{} (CDIAL 4024).
\end{cousinweb}

\begin{grammarlens}[title={What you've built}]
One more word for this chapter.
\end{grammarlens}

\subsection*{Your turn}
\begin{itemize}
\raggedright
  \item \emph{Say it:} \emph{gahrā}
  \item \emph{Say it:} \emph{gahrā}, once more
  \item \emph{Say it:} say \emph{gahrā}
  \item \emph{Recall:} say the Urdu for wide, then the Urdu for narrow, then say \emph{gahrā} again
\end{itemize}

\begin{tcolorbox}[breakable,colback=teal!4,colframe=teal!35!black,title={Before you move on}]
What does \ur{گہرا} mean? (\textquotedblleft{}Deep, dark (of colour)\textquotedblright{}.) Say it once more.
\end{tcolorbox}
